% !PRINTED 55
\weberpage{55}
%<<CONT

then carrying out the multiplication of the two formulae (4), and seeking in the product the coefficient of $y^\nu$, gives
\begin{equation}
 {D_\nu(f\varphi)\over\Pi(\nu)}
 =u_\nu v_0+u_{\nu-1}v_1+u_{\nu-2}v_2+\cdots+u_0v_\nu,
\tag{6}
\end{equation}
or
\begin{equation}
 D_\nu(f\varphi)=\varphi D_\nu f+B_1^{(\nu)}D_{\nu-1}fD_1\varphi
 +B_2^{(\nu)}D_{\nu-2}fD_2\varphi+\cdots,
\tag{7}
\end{equation}
where $B_1^{(\nu)},B_2^{(\nu)},\ldots$ are the binomial coefficients. In a similar way one can, by using the polynomial coefficients, form the derivatives of a product of more than two factors.

We shall form the \gesperrt{first derivative}, which we now denote by $D$ instead of by $D_1$, for a product of $n$ factors. For two factors we obtain from (7)
\[
 D(f\varphi)=\varphi Df+fD\varphi=\varphi f'(x)+f\varphi'(x),
\]
and generally, if we denote the factors by $u_1,u_2,\ldots,u_n$ and their derivatives by $u'_1,u'_2,\ldots,u'_n$,
\begin{equation}
\begin{aligned}
D(u_1u_2\cdots u_n)
&=u'_1u_2\cdots u_n+u_1u'_2\cdots u_n+\cdots+u_1u_2\cdots u'_n,
\end{aligned}
\tag{8}
\end{equation}
or more briefly,
\[
 {D(u_1u_2\cdots u_n)\over u_1u_2\cdots u_n}
 =\sum {D u_\nu\over u_\nu}.
\]

Thus if we consider a product of linear factors
\begin{equation}
 f(x)=(x-\alpha_1)(x-\alpha_2)\cdots(x-\alpha_n),
\tag{9}
\end{equation}
we obtain, since the first derivatives of $x-\alpha_1$, $x-\alpha_2,\ldots,x-\alpha_n$ are all equal to $1$,
\begin{equation}
\begin{aligned}
f'(x)&=(x-\alpha_2)(x-\alpha_3)\cdots(x-\alpha_n)\\
&\quad +(x-\alpha_1)(x-\alpha_3)\cdots(x-\alpha_n)+\cdots\\
&\quad +(x-\alpha_1)(x-\alpha_2)\cdots(x-\alpha_{n-1}),
\end{aligned}
\tag{10}
\end{equation}
for which one can also write
\begin{equation}
 f'(x)={f(x)\over x-\alpha_1}+{f(x)\over x-\alpha_2}+\cdots+{f(x)\over x-\alpha_n}.
\tag{11}
\end{equation}
A very important result follows from this if $x$ is put equal to one of the values $\alpha_1,\alpha_2,\ldots,\alpha_n$, namely
%CONT>>
